\section{Vendor Mitigation and Comparative Evaluation}
\label{sec:vendor-mitigation}

Chapter~7 showed that a custom compensating control could interrupt the tested
Copy Fail paths while the vulnerable kernel remained active. It also showed why
the word \emph{custom} matters. MORI Guard v1 achieved broad prevention by
removing access to \code{algif\_aead}, whereas MORI v2.7 preserved the tested
interface operations by enforcing a narrower behavioural correlation. RQ5 now
places those results beside Canonical's temporary response:

\begin{quote}
\textbf{RQ5 -- Vendor Comparison.} How does the custom compensating control
compare with Canonical's temporary mitigation, which prevents
\code{algif\_aead} from being loaded?
\end{quote}

The comparison is not a contest in which the most selective mechanism wins by
definition. A broad vendor mitigation may be preferable because it is simple,
packaged, supportable, and intentionally conservative. A selective research
control may preserve more functionality, but only within the behaviour and
environment that its policy actually recognises. The evaluation must therefore
consider prevention scope, compatibility, operational complexity, evidential
strength, telemetry, and the route back to an ordinary patched system.

The resulting answer is deliberately asymmetric. MORI Guard v1 belongs to the
same policy class as Canonical's temporary mitigation and offers no compatibility
advantage over it. MORI v2.7 demonstrates that the tested exploitation path can
be interrupted without removing the tested AF\_ALG AEAD operations, but it does
so with narrower assurance and substantially greater implementation and
deployment complexity. Neither custom control corrects the kernel defect, and
neither should displace the vendor kernel update as the final remediation.

\subsection{Comparison Objective and Evidence Boundary}

The comparison uses three evidence classes. First, Canonical's security notice
and technical guidance define the purpose, package version, deployment action,
compatibility warning, and intended lifetime of the temporary mitigation
\cite{ubuntu-copyfail-fixes,ubuntu-copyfail-cve}. Second, Phase~01 records the
policy as it appeared on the preserved Ubuntu reference guest. Third, Phase~04
contains the direct vulnerable-kernel experiments for MORI Guard v1 and MORI
v2.7.

These classes answer different questions. Vendor documentation states what the
distribution intended to deploy. Phase~01 shows the effective module-resolution
rule on the reference system. Phase~04 shows what happened when the laboratory's
custom controls were exercised on \code{6.8.0-116-generic}. None of these alone
creates a perfectly matched head-to-head trial.

In particular, the Canonical \code{kmod} policy was captured on a reference
guest whose running \code{6.8.0-137-generic} kernel was already corrected. A
Copy Fail execution on that combined state would have had two independent
reasons to fail: the ordinary module-loading path was restricted and the kernel
contained the vendor correction. The project therefore did not use that guest
as a causal exploit test for the temporary mitigation. By contrast, MORI Guard
v1 was directly tested on \code{6.8.0-116-generic}, and the Phase~04 chronology
places the final MORI v2.7 validation before vendor remediation. The final v2.7
captures identify the frozen binary and result but do not repeat \code{uname};
their kernel attribution is therefore sequence-based rather than a repeated
runtime capture.

The comparison criteria are summarised in \cref{tab:vendor-comparison-criteria}.

\begin{table}[H]
    \centering
    \caption{Criteria used for the vendor-mitigation comparison.}
    \label{tab:vendor-comparison-criteria}
    \small
    \begin{tabularx}{\textwidth}{@{}>{\bfseries\raggedright\arraybackslash}p{0.23\textwidth}>{\raggedright\arraybackslash}p{0.32\textwidth}Y@{}}
        \toprule
        Criterion & Question & Boundary \\
        \midrule
        Prevention scope &
        Which prerequisite or behaviour is denied? &
        Broader coverage can reduce evasion space while increasing collateral
        impact. \\

        Compatibility &
        Which benign operations remain available? &
        Only operations actually exercised by the preserved controls are
        credited. \\

        Evidential strength &
        Was the result directly tested on the vulnerable runtime or inferred
        from the mechanism and vendor record? &
        Mechanism-based confidence is kept separate from experimental causality. \\

        Operational burden &
        What must be installed, loaded, maintained, and recovered? &
        Laboratory feasibility does not establish production suitability. \\

        Observability &
        What security-relevant context accompanies a denial? &
        Telemetry is useful, but it does not compensate for an incomplete
        prevention boundary. \\

        Remediation lifecycle &
        How is the control retired after a fixed kernel is active? &
        A temporary control should have a clear and supportable exit path. \\
        \bottomrule
    \end{tabularx}
\end{table}

This boundary prevents two overstatements. The chapter does not claim that the
Canonical rule was independently proven as the sole cause of exploit failure
on the vulnerable laboratory kernel. It also does not treat the successful
MORI v2.7 controls as proof that every legitimate AF\_ALG workload remains
compatible or that every Copy Fail implementation is recognised.

\subsection{Canonical's Temporary \texttt{kmod} Mitigation}

Canonical published \code{kmod} updates on 30 April 2026 to block loading of
the affected \code{algif\_aead} module until kernel updates became available.
For Ubuntu 24.04 LTS Noble, the mitigation-bearing package version was
\code{31+20240202-2ubuntu7.2} \cite{ubuntu-copyfail-fixes}. Canonical's manual
equivalent writes the following \code{modprobe} install rule:

\begin{lstlisting}
install algif_aead /bin/false
\end{lstlisting}

An install rule substitutes the specified command for the module insertion
performed through normal \code{modprobe} resolution. Requests may still resolve
the dependency path, but the \code{algif\_aead} insertion itself is replaced by
\code{/bin/false}. The mechanism does not inspect the caller, target file,
splice activity, or intended cryptographic operation. It removes the affected
module from that ordinary loading path for every caller.

The rule is preventive only when the module is not already active. Canonical's
guidance therefore instructs administrators to unload \code{algif\_aead} when
possible or reboot when the module is in use. The same guidance warns that
applications may fall back to userspace cryptography, but that the fallback is
not guaranteed and already running applications may be affected
\cite{ubuntu-copyfail-fixes}. This is not an incidental weakness in the
mitigation; it is the explicit trade-off of removing an attack surface before a
corrected kernel can be deployed.

\subsubsection{Observed reference-system policy}

Phase~01 captured the following state on the Ubuntu reference guest:

\begin{lstlisting}
kmod                                  31+20240202-2ubuntu7.2
linux-image-6.8.0-137-generic         6.8.0-137.137

/etc/modprobe.d/disable-algif_aead.conf:
install algif_aead /bin/false
\end{lstlisting}

The corresponding dry run resolved to insertion of the dependency
\code{af\_alg} followed by \code{install /bin/false}; the module-state capture
contained no loaded \code{algif\_aead} entry. The installed kernel configuration
also recorded \code{CONFIG\_CRYPTO\_USER\_API\_AEAD=m}, so the affected
interface was a loadable module on this system rather than compiled directly
into the kernel.

This evidence confirms the effective policy and its suitability for the
captured Ubuntu configuration. It does not establish resistance to an
administrator deliberately changing the configuration, nor does it test every
low-level way privileged code might be introduced. Those cases fall outside the
unprivileged local-attacker model used by the project. More importantly for RQ5,
the capture proves policy presence and resolution behaviour, not a standalone
vulnerable-kernel exploit result.

\subsubsection{Intended lifetime}

Canonical described the \code{kmod} change as temporary and recommended
upgrading to a fixed kernel and rebooting. Its published Noble generic-kernel
boundary identifies \code{6.8.0-117.117} as fixed
\cite{ubuntu-copyfail-cve}. Once a corrected kernel is active, the loading rule
is no longer required and may itself become an unnecessary compatibility
restriction. Canonical's guidance consequently includes a documented route for
disabling the mitigation after the kernel update
\cite{ubuntu-copyfail-fixes}.

This lifecycle is a significant operational advantage. The mitigation is
distributed through the ordinary package channel, tied to a vendor advisory,
and explicitly superseded by a kernel update. Its purpose is emergency risk
reduction, not permanent replacement of AF\_ALG AEAD.

\subsection{Relationship to MORI Guard v1}

MORI Guard v1 used the same system mechanism but redirected the module request
to a laboratory blocker:

\begin{lstlisting}
install algif_aead /usr/local/sbin/mori-block-algif-aead
\end{lstlisting}

The blocker returned failure, wrote a structured journal event, and attempted
to present the MORI message to the requesting process. On
\code{6.8.0-116-generic}, the Python proof of concept failed, \code{labuser}
remained UID~1001, and \code{/usr/bin/su} retained its known-good digest. A
separate unprivileged \code{gcm(aes)} bind also failed while the guard was
active.

At the policy level, Canonical's mitigation and MORI Guard v1 are therefore
equivalent in breadth for the tested Ubuntu module configuration: both replace
ordinary loading of \code{algif\_aead} with a failing command. The custom script
changes the presentation and evidence available to the laboratory, not the
fundamental compatibility boundary. Both deny malicious and benign callers
before either can use the module.

There are nevertheless operational differences:

\begin{itemize}
    \item Canonical's rule arrives in a signed distribution package and is
    documented as part of the vendor response. MORI Guard v1 is a locally
    authored configuration and script whose review, deployment, and removal
    remain the operator's responsibility.

    \item MORI Guard v1 emits a dedicated journal record and a recognisable
    presentation message. Canonical's rule provides a generic command failure;
    the vendor guidance does not define equivalent structured Copy Fail deny
    telemetry.

    \item MORI Guard v1 was directly exercised against the vulnerable kernel in
    the laboratory. The captured Canonical rule was validated on the already
    patched reference guest and is supported by the vendor's stated mechanism,
    not by an otherwise identical laboratory exploit trial.

    \item Both depend on the module being absent when enforcement begins. A
    loading rule does not retroactively detach an active module, so operational
    validation must include module state or a reboot.
\end{itemize}

For an ordinary Ubuntu deployment awaiting a kernel update, MORI Guard v1 does
not offer a security or compatibility reason to replace the vendor package. Its
additional value is limited to the research environment, where the custom
blocker made the control decision visible and provided a locally controlled
comparison baseline.

\subsection{Relationship to MORI v2.7}

MORI v2.7 begins from the opposite compatibility decision. It leaves
\code{algif\_aead} available and denies only when the tested same-process
correlation reaches the protected privileged-file read. The frozen candidate
combines live AEAD socket state, exact-invocation SUID-file-to-pipe splice
state, a non-root read of the selected class of root-owned SUID file, and an
authenticated monitor exemption.

This design produced two useful differences from the module-level controls:

\begin{enumerate}
    \item The unprivileged \code{gcm(aes)} socket, bind, and later recorded
    accept operations remained available where tested. AEAD-only and
    splice-only negative controls did not reach the deny decision.

    \item A matched decision produced \code{EPERM} together with kernel trace
    output and a structured ring-buffer event containing the TGID, UID,
    timestamp, and deny type. The telemetry exposed why the policy intervened
    rather than only reporting that module insertion failed.
\end{enumerate}

The selectivity is real within the preserved test matrix, but it changes the
assurance model. Canonical's mitigation prevents the complete module-dependent
path through ordinary module resolution. MORI v2.7 must correctly observe and
correlate the exploit's behaviour. A variant that separates work across
processes, changes ownership semantics, avoids a monitored hook, or falls
outside the selected SUID-to-pipe sequence may evade the policy. Conversely, a
legitimate program that reproduces the full correlation may receive
\code{EPERM}.

MORI v2.7 also introduces dependencies that the \code{kmod} rule does not have:
BPF LSM support, BTF and CO-RE information, the selected LSM and tracing hooks,
a privileged userspace loader, map allocation, state lifecycle correctness,
and optional monitor authentication. The Phase~04 source review further shows
that state-map update failures preserve the original operation. That choice
avoids turning allocation pressure into an uncontrolled denial of service, but
it is fail-open for the compensating policy.

The appropriate conclusion is therefore not that selective correlation is
strictly better than broad removal. MORI v2.7 trades some breadth and
operational simplicity for retained tested functionality and richer decision
context. Whether that exchange is acceptable depends on the system's workload,
kernel facilities, risk tolerance, and the length of time for which the
vulnerable kernel must remain in service.

\subsection{Security Coverage and Failure Modes}

The three controls can be ordered by where they intervene, but not by a single
universal measure of strength. \Cref{tab:vendor-security-comparison} separates
the supported security claims.

\begin{table}[H]
    \centering
    \caption{Security boundary of the temporary controls.}
    \label{tab:vendor-security-comparison}
    \scriptsize
    \begin{tabularx}{\textwidth}{@{}>{\bfseries\raggedright\arraybackslash}p{0.16\textwidth}>{\raggedright\arraybackslash}p{0.25\textwidth}>{\raggedright\arraybackslash}p{0.25\textwidth}Y@{}}
        \toprule
        Dimension & Canonical \code{kmod} mitigation & MORI Guard v1 & MORI v2.7 \\
        \midrule
        Intervention point &
        Normal \code{modprobe} resolution for \code{algif\_aead}. &
        Same resolution path, redirected to a custom blocker. &
        Correlated privileged-file read inside a qualifying splice operation. \\

        Protected scope &
        All callers that require the blocked module through the captured loading
        path. &
        Same broad module dependency for the tested laboratory path. &
        Tested same-process AEAD, SUID-to-pipe splice, and privileged-read
        sequence. \\

        Vulnerable code removed &
        No. Access is restricted while the vulnerable kernel may remain. &
        No. The vulnerable kernel is deliberately retained. &
        No. The policy denies selected behaviour before the target changes. \\

        Direct vulnerable-kernel PoC evidence &
        Not isolated as the sole variable in this laboratory. &
        Python path blocked on \code{6.8.0-116-generic}. &
        Python and independent C paths denied in Phase~04 before vendor
        remediation; the final captures do not repeat \code{uname}. \\

        Principal bypass boundary &
        Module already active or loading outside the validated ordinary path. &
        Same, plus local script and configuration integrity. &
        Behaviour outside the selected hooks, state, process ownership, or
        lifecycle model. \\

        False-positive surface &
        Any benign caller requiring the module. &
        Any benign caller requiring the module. &
        A benign caller reproducing the full correlated sequence. \\
        \bottomrule
    \end{tabularx}
\end{table}

For the public proof of concept's known dependency, module-level removal has a
simple advantage: the exploit cannot reach \code{algif\_aead} through the
validated loading route. There is no correlation window, per-process state, or
target predicate to evade. The cost is equally simple: neither can a benign
caller reach that module through the same route.

MORI v2.7 narrows both the denial and the supported prevention statement. Its
two final proof-of-concept results, negative controls, target sweep, and retained
lifecycle regression increase confidence in the tested path. They do not equal
the coverage obtained by removing the prerequisite interface. This distinction
matters most during emergency response, where predictable containment may be
more valuable than preserving a specialised interface for a short period.

\subsection{Compatibility, Operations, and Telemetry}

Compatibility evidence is intentionally graded rather than binary. Canonical
warned that applications should fall back to userspace cryptography but may not
implement that fallback. The project did not infer successful application
behaviour from the word \emph{should}. MORI Guard v1 directly demonstrated the
loss of one benign unprivileged \code{gcm(aes)} bind. MORI v2.7 preserved the
tested socket, bind, and accept operations, but the preserved Phase~04 controls
do not contain a complete encryption and decryption round trip or a
representative application workload.

The result is a bounded compatibility ordering:

\begin{lstlisting}
Canonical kmod mitigation / MORI Guard v1
    tested interface removed through normal module resolution

MORI v2.7
    tested socket, bind, and accept operations retained
    complete application compatibility not established
\end{lstlisting}

Operational characteristics are compared in
\cref{tab:vendor-operational-comparison}.

\begin{table}[H]
    \centering
    \caption{Operational comparison of the temporary controls.}
    \label{tab:vendor-operational-comparison}
    \scriptsize
    \begin{tabularx}{\textwidth}{@{}>{\bfseries\raggedright\arraybackslash}p{0.16\textwidth}>{\raggedright\arraybackslash}p{0.25\textwidth}>{\raggedright\arraybackslash}p{0.25\textwidth}Y@{}}
        \toprule
        Dimension & Canonical \code{kmod} mitigation & MORI Guard v1 & MORI v2.7 \\
        \midrule
        Delivery &
        Ubuntu package and vendor advisory. &
        Local configuration plus root-owned blocker script. &
        Local BPF build, privileged loader, service integration, and policy
        state. \\

        Initial activation &
        Package update plus unload or reboot where required. &
        Deploy rule and script, confirm module absence, then validate. &
        Verify kernel facilities and monitor trust, load BPF object, attach
        programs, and keep loader active. \\

        Decision telemetry &
        Generic module-load failure; no dedicated structured deny event is
        defined by the captured rule. &
        Journal event and optional presentation message. &
        Trace event, structured ring-buffer record, and optional TTY message. \\

        Maintenance &
        Vendor tracks package lifecycle and superseding kernel update. &
        Operator owns script, configuration, and local rollback. &
        Operator owns source, build chain, hooks, maps, trust model, regressions,
        and portability. \\

        Failure tendency &
        Conservative and broad: unavailable interface. &
        Conservative and broad, with an additional local script dependency. &
        Selective but potentially fail-open on state-allocation failure or
        out-of-model behaviour. \\

        Retirement &
        Documented after installation of a fixed kernel. &
        Remove local rule and blocker after patch validation. &
        Detach programs, stop loader, remove service state, and repeat benign
        and adversarial validation on the patched kernel. \\
        \bottomrule
    \end{tabularx}
\end{table}

MORI's richer telemetry is useful for research and incident reconstruction. It
can distinguish a matched behavioural policy from a generic module-loading
failure and can preserve the identity associated with a denial. That advantage
does not make it a stronger emergency mitigation by itself. A simpler control
with less context may still be preferable when its broader denial boundary is
acceptable and its deployment path is substantially easier to support.

The same reasoning applies to vendor status. Packaging and vendor support do not
prove that a control has no defects, but they materially change who maintains
the response, how updates are distributed, and how an administrator determines
when the workaround should be removed. MORI v2.7 has no equivalent portability,
performance, or long-term maintenance programme.

\subsection{Decision Analysis}

The comparison produces three practical decisions rather than one unconditional
winner.

\subsubsection{Emergency containment before a patch}

Where applications do not require the affected interface, or where temporary
loss is acceptable, Canonical's packaged mitigation is the preferred emergency
control. It removes the known module prerequisite through a standard Ubuntu
update and has a documented activation and retirement process. MORI Guard v1
does not improve that compatibility boundary and adds locally maintained code.

\subsubsection{Temporary retention of tested interface operations}

Where the vulnerable kernel must temporarily remain and the tested AF\_ALG AEAD
operations cannot be removed, MORI v2.7 demonstrates a narrower experimental
option. The decision should be conditioned on the actual workload matching the
tested environment, acceptance of the identified bypass and false-positive
boundaries, availability of the required kernel facilities, and active
ownership of the control's maintenance.

This is a research result, not a general deployment recommendation. A production
operator would need additional application compatibility testing, sustained
performance measurement, concurrency and lifecycle regression, fail-open and
fail-closed analysis, independent code review, and a supported deployment and
rollback process.

\subsubsection{Long-term remediation}

Neither form of temporary denial should remain the target state. The correct
long-term action is installation of a fixed vendor kernel and removal of the
workaround after reboot and validation. This restores ordinary interface
availability while addressing the in-place AEAD behaviour in the kernel rather
than predicting or preventing one exploitation sequence.

\subsection{Answer to RQ5}

RQ5 asks how the custom compensating control compares with Canonical's temporary
module-loading mitigation. The answer depends on which MORI generation is being
compared.

MORI Guard v1 is functionally similar to Canonical's mitigation for the tested
Ubuntu module configuration. Both prevent normal loading of
\code{algif\_aead}, interrupt the known module-dependent exploitation path, and
remove benign access to the same interface. MORI Guard v1 adds laboratory
journalling and presentation, but it is locally maintained and provides no
compatibility advantage. Canonical's packaged rule is therefore the stronger
operational choice for broad temporary containment.

MORI v2.7 offers a different trade-off. It preserved the tested AF\_ALG AEAD
socket, bind, and accept operations and denied the correlated Python and
independent C proof-of-concept paths in the pre-remediation Phase~04 sequence.
It also supplied structured decision telemetry. Those advantages are balanced by a narrower
behavioural coverage claim, more complex privileged deployment, kernel-feature
dependencies, possible false positives for a benign matching sequence, and
explicit evasion and fail-open boundaries.

The supported conclusion is therefore:

\begin{lstlisting}
Canonical's kmod mitigation is the preferable vendor-supported
emergency containment when broad module removal is acceptable.

MORI v2.7 demonstrates that narrower prevention with greater
tested interface availability is possible, but only as an
experimental compensating control with bounded assurance.

Neither custom control replaces the fixed vendor kernel.
\end{lstlisting}

The laboratory does not claim that MORI v2.7 is universally safer, more secure,
or more compatible than Canonical's response. It demonstrates a defensible
engineering alternative for the tested path and makes the cost of selectivity
visible.

\subsection{Transition to Patched-Kernel Validation}

RQ5 ends at the temporary-control decision. RQ6 asks a different question:
whether the reproduced behaviour disappears after installation of a corrected
kernel and whether the final observations match the expected upstream change.
Canonical identifies \code{6.8.0-117.117} as the first fixed Noble generic
kernel; the Phase~05 guest ran the later \code{6.8.0-137.137} package
\cite{ubuntu-copyfail-cve}.

That transition must isolate the kernel correction from both custom and
module-level denial. The patched retest therefore needs to show that MORI is not
making the decision, that the Canonical loading rule is disabled for the test,
that \code{algif\_aead} can load, and that the benign AF\_ALG control reaches
the operations exercised by the laboratory. Only then can a failed proof of
concept be attributed to the patched execution path rather than to continued
interface removal.

The Phase~05 package also contains three provenance imperfections that Chapter~9
must preserve rather than silently repair:

\begin{itemize}
    \item The file named \code{phase05-pre-patch.txt} already records
    \code{6.8.0-137-generic}. It is a patched pre-retest baseline, not evidence
    of a pre-patch runtime.

    \item The machine-readable C provenance capture contains permission errors
    instead of the claimed source and binary hashes. The successful values are
    preserved in screenshot~10 and must be attributed to that visual evidence.

    \item The timestamped machine-readable MORI-isolation capture was produced
    later than the timestamped C provenance capture. Screenshot~02 records the
    isolation state but has no embedded terminal time, so strict pre-test causal
    ordering is not established by timestamps alone.
\end{itemize}

These caveats do not change the RQ5 comparison, which rests on the Phase~01
vendor-policy capture, the Phase~04 control experiments, and Canonical's
published response. They do constrain the wording of the final patch-validation
claim. Chapter~9 therefore treats the Phase~05 result as a bounded empirical
retest, distinguishes successful screenshots from failed text captures, and
separates observed state from inferred chronology.